\documentclass{fiche}

\surtitre{Fiche 3 · Proportionnalité et fonctions}
\titrefiche{La proportionnalité}
\accroche{Si une seule leçon devait être maîtrisée dans toute l'année, c'est
celle-ci. Recettes, prix, échelles, vitesses, pourcentages, statistiques : elle
est partout, au collège comme dans la vie courante.}
\niveaux{6\ieme{} · 5\ieme{}}
\priorite{3}
\pourquoi{La notion la plus rentable du programme : elle irrigue tous les autres chapitres.}
\duree{35 min}
\domaine{Proportionnalité, fonctions}

\begin{document}
\entetefiche

\begin{automatismes}[2]
  \auto[6\ieme]{Repérer double, moitié, tiers, quart entre deux nombres}
  \auto[6\ieme]{« 4 fois plus » \fleche{} multiplier ; « 4 fois moins » \fleche{} diviser}
  \auto[5\ieme]{Dire si une situation relève ou non de la proportionnalité}
  \auto[5\ieme]{Recette pour 4 personnes \fleche{} la donner pour 2, 6 ou 8}
  \auto[5\ieme]{Prix d'un kilo connu \fleche{} prix de 3 kg, de 4{,}3 kg}
  \auto[5\ieme]{Prendre 1 \%, 10 \% ou 50 \% d'un nombre}
  \auto[règle]{Le coefficient se lit toujours « par unité » : prix \emph{au} kilo,
               battements \emph{par} minute}
\end{automatismes}

\section{Reconnaître une situation de proportionnalité}

\begin{definition}
Deux grandeurs sont \motcle{proportionnelles} lorsque l'on passe de l'une à
l'autre en multipliant \emph{toujours par le même nombre}. Ce nombre s'appelle
le \motcle{coefficient de proportionnalité}.
\end{definition}

\begin{tableaufiche}{@{}G{58mm} Y@{}}
\ligneentete \entetecell{C'est proportionnel} & \entetecell{Ce ne l'est pas}\\
Le prix payé et la quantité achetée & L'âge et la taille d'une personne\\
La distance parcourue et la durée, à vitesse constante & Le prix d'un abonnement avec frais d'inscription\\
Le périmètre d'un carré et son côté & L'aire d'un carré et son côté\\
\end{tableaufiche}

\begin{piege}[Le test qui ne trompe pas]
Si l'on double la première grandeur, la seconde \emph{doit} doubler.
Un carré de côté 2 a pour aire 4 ; de côté 4, son aire est 16, et non 8 :
l'aire n'est \textbf{pas} proportionnelle au côté.
\end{piege}

\section{Les trois procédures du programme}

Le programme demande de \emph{choisir la procédure adaptée aux nombres},
et non d'en appliquer une seule mécaniquement.

\subsection{Linéarité multiplicative — « 3 fois plus »}
5 croissants coûtent 6{,}25 €. Combien coûtent 15 croissants ?
\[15 = 5\times 3 \quad\text{donc}\quad 6{,}25\times 3 = 18{,}75\ \text{€}\]

\subsection{Linéarité additive — « 5 + 3 »}
5 croissants : 6{,}25 €. 3 croissants : 3{,}75 €. Alors 8 croissants :
\[6{,}25 + 3{,}75 = 10\ \text{€}\]

\subsection{Retour à l'unité — le passe-partout}
\begin{methode}
\begin{enumerate}[leftmargin=6mm,label=\textbf{\arabic*.}]
  \item Chercher la valeur pour \textbf{une} unité.
  \item Multiplier par le nombre voulu.
\end{enumerate}
\vspace{0.8mm}
5 croissants coûtent 6{,}25 € \fleche{} un croissant coûte
$6{,}25\div 5 = 1{,}25$ € \fleche{} 8 croissants coûtent
$1{,}25\times 8 = 10$ €.
\end{methode}

\begin{center}
\tableaupropo{Croissants}{Prix (€)}{5/6{,}25, 8/?, 12/15}{$\times\,1{,}25$}
\end{center}

\begin{piege}[Ce que le programme interdit explicitement en 6\ieme{}]
La technique du « produit en croix » \textbf{n'est pas enseignée} : le texte
officiel demande que l'élève raisonne sur le \emph{sens} de la situation.
Un résultat juste obtenu par produit en croix ne vaut donc pas la démarche attendue.
\end{piege}

\section{Représenter une situation}

\subsection{Par un tableau}
Le nom de chaque grandeur, avec son unité, figure explicitement dans le tableau.
Le programme y insiste : c'est ce qui donne du sens aux nombres.

\subsection{\niv{5\ieme} Par un graphique}

\begin{retenir}[La signature graphique de la proportionnalité]
Les points sont \textbf{alignés} \emph{et} l'alignement \textbf{passe par l'origine}.
Les deux conditions sont nécessaires : des points alignés sur une droite qui ne
passe pas par l'origine ne relèvent pas de la proportionnalité.
\end{retenir}

\begin{center}
\repereplan[7mm]{6}{5}{1/1/, 2/2/, 3/3/, 4/4/}
\end{center}

\section{Les grandes applications}

\subsection{Les pourcentages}
Appliquer un pourcentage, c'est utiliser un coefficient :
\[15\,\%\ \text{de }240 \;=\; \frac{15}{100}\times 240 \;=\; 36\]
Voir la fiche 2 pour le détail.

\subsection{Les échelles}
\begin{definition}
Une échelle est le coefficient qui fait passer de la réalité au dessin :
\[\text{échelle} = \frac{\text{distance sur le plan}}{\text{distance réelle}}
\quad\text{(dans la même unité)}\]
\end{definition}

À l'échelle $\tfrac{1}{25\,000}$, 1 cm sur la carte représente 25 000 cm dans la
réalité, soit \textbf{250 m}. Donc 4 cm sur la carte représentent \textbf{1 km}.

\begin{piege}
L'erreur classique est d'oublier la conversion : 25 000 cm ne font pas 25 000 m.
On convertit \emph{toujours} avant de conclure.
\end{piege}

\subsection{\niv{5\ieme} La vitesse moyenne}
La vitesse est le coefficient entre la distance et la durée :
\[v = \frac{d}{t} \qquad d = v\times t \qquad t = \frac{d}{v}\]
180 km parcourus en 2 h 30 : attention, 2 h 30 vaut $2{,}5$ h, et non $2{,}30$ h.
\[v = \frac{180}{2{,}5} = 72\ \text{km/h}\]

\begin{piege}
2 h 30 s'écrit $2{,}5$ h en écriture décimale : les durées sont en base 60,
pas en base 10. C'est l'une des erreurs les plus coûteuses du collège.
\end{piege}

\section{\niv{5\ieme} L'expression « en fonction de »}

En 5\ieme{}, on commence à dire qu'une grandeur \motcle{dépend} d'une autre :
« le prix \emph{en fonction de} la quantité ». On produit alors une formule simple,
un tableau de valeurs, ou un graphique — trois façons de dire la même chose.

\[\text{prix} = 1{,}25 \times \text{nombre de croissants}\]

La proportionnalité devient ainsi le premier exemple de fonction, notion qui sera
étudiée en 4\ieme{} et en 3\ieme{}.

\begin{prolongement}
Le programme relie la proportionnalité aux grands enjeux contemporains :
lire un graphique d'émissions de CO\textsubscript{2}, comparer des ordres de
grandeur, repérer une présentation trompeuse. C'est là que les mathématiques
deviennent un outil d'esprit critique.
\end{prolongement}

\begin{videos}
  \video{QgjbpX_kciA}{Reconnaître la proportionnalité — 6\ieme}{Yvan Monka}{7:19}{275 000}
  \video{u4lVY1tGSnA}{Proportionnalité — 6\ieme}{Les Bons Profs}{4:06}{264 000}
  \video{Qy2ppBOEax4}{Appliquer une proportionnalité — 6\ieme}{Yvan Monka}{6:30}{112 000}
\end{videos}

\end{document}
